% Working draft: the revised readout and warm-up are proposed, not yet tested.
% Results are provisional working material for the latent-only control (PGND-0).
% Select the legacy IEEE Times encoding before the class initializes its fonts.
\RequirePackage[T1]{fontenc}
\documentclass[journal,twoside]{IEEEtran}
\usepackage{amsmath,amssymb,bm}
\usepackage[noadjust]{cite}
\usepackage[ruled,vlined]{algorithm2e}
\usepackage[hidelinks]{hyperref}
\usepackage{xurl}

\title{Physics-Guided Neural Dynamics for Pantograph--Catenary Contact Force Estimation}

\author{Laha Ale,~\IEEEmembership{Senior~Member,~IEEE},
Yang Song,~\IEEEmembership{Senior~Member,~IEEE},
and Zhigang Liu,~\IEEEmembership{Life~Fellow,~IEEE}
\thanks{L. Ale is with the School of Computing and Artificial Intelligence,
Southwest Jiaotong University, Chengdu, China (e-mail: laha.ale@ieee.org).}
\thanks{Y. Song and Z. Liu are with the School of Electrical Engineering,
Southwest Jiaotong University, Chengdu, China (e-mail:
y.song\_ac@hotmail.com; liuzg\_cd@126.com).}
\thanks{L. Ale and Y. Song are with the SWJTU-Leeds Joint School,
Southwest Jiaotong University, Chengdu, China (e-mail:
laha.ale@ieee.org; y.song\_ac@hotmail.com).}
\thanks{Corresponding author: Yang Song (e-mail: y.song\_ac@hotmail.com).}}

\begin{document}
\maketitle

\begin{abstract}
Estimating pantograph--catenary contact force from panhead acceleration and
uplift offers an alternative to direct force instrumentation. This work
investigates physics-guided neural dynamics (PGND) using simulation data.
The model evolves a latent configuration--velocity state with
positive-semidefinite damping, an energy-derived restoring term,
observation-conditioned driving, and a regularized neural residual.
We propose an observation-conditioned readout combining a direct prediction
from the latest available response with a correction decoded from the
latent dynamics. The information budget is restricted to a finite
observation history when predicting the next force sample. A causal warm-up
initializer is formulated as a separate extension. The evaluation protocol
specifies common recording-level test partitions, training-only scaling,
and force-error metrics for PGND and recurrent baselines. Controlled
ablations are designed to distinguish the effects of observation access,
temporal context, and mechanical structure. Quantitative evaluation of the
revised method remains pending; the present draft makes no claim of
improved accuracy or robustness.
\end{abstract}

\begin{IEEEkeywords}
Pantograph--catenary interaction, contact force estimation,
physics-guided learning, neural ordinary differential equations,
simulation-based evaluation.
\end{IEEEkeywords}
\IEEEpeerreviewmaketitle
\thispagestyle{plain}
\pagestyle{plain}

\section{Introduction}
\label{sec_intro}

The pantograph--catenary system transfers electrical power from the overhead
line to a moving train. Its contact-force history is an important indicator
of current collection quality and mechanical interaction
\cite{liu2024review,bruni2018interaction,navik2017variation}.
Excessive force fluctuations accelerate wear, whereas insufficient force
can lead to contact loss and arcing \cite{bruni2018interaction}.
Contact interaction is therefore considered in technical assessment
procedures such as EN~50367 \cite{en50367_2020}.

Direct force measurement uses an instrumented pantograph and accounts for
support forces and inertial and aerodynamic contributions
\cite{en50317_2012}. Panhead acceleration and uplift provide alternative
response measurements containing information about the coupled dynamics
\cite{carnevale2016processing,blanco2022panhead}. Learning this relationship
can support indirect force estimation; related work has also used
acceleration to estimate contact-force statistics
\cite{gregori2023assessment}.

Our previous study investigated RNN, LSTM, GRU, and CNN--GRU predictors using
responses generated by a numerical pantograph--catenary model
\cite{ale2026realtime}. That study considered different catenary
configurations, speeds, and filtering bandwidths. The present work reuses
the available simulation files but adopts a documented, shared preprocessing
and evaluation protocol. Because file handling, window boundaries, and
normalization differ from the legacy implementation, the new scores are
not presented as a reproduction of the published scores.

Discrete sequence models can learn accurate response--force mappings
without explicitly representing mechanical evolution. Neural ordinary
differential equations (ODEs) instead parameterize a continuous vector field
\cite{chen2018neuralode}. Their continuous-time representation does not,
by itself, ensure physical consistency or superior prediction. Available
mechanical knowledge can guide the vector field, following the broader
principle of combining scientific structure with learning
\cite{karniadakis2021physics}. The second-order mechanics suggest a
configuration--velocity state, dissipative terms, and energy-derived
restoring terms \cite{bruni2015benchmark,bruni2018interaction}.
These provide an inductive bias, not an identified reduced-order model of
the full pantograph--catenary system.

The original PGND implementation routes all observation information through
the latent dynamics before decoding force. Our preliminary comparison does
not demonstrate an accuracy advantage over CNN--GRU. This motivates the
testable hypothesis that a direct observation-conditioned prediction path
can complement the latent dynamical representation. We therefore propose
an additive readout and separately formulate initialization from a short
available history. Neither modification is assumed to improve accuracy:
direct-only and unstructured-ODE ablations are needed to determine whether
any gain is attributable to mechanical structure.

The contributions and present scope of this draft are as follows:
\begin{itemize}
    \item A finite-history PGND formulation combines structured latent
    dynamics with an observation-conditioned readout while excluding
    target-time observations.
    \item A causal warm-up initializer and standardized-force objective
    are specified, with explicit limits on the physical interpretation of
    latent energy and the two output branches.
    \item A common simulation-data protocol and measured results for the
    original formulation are reported. A staged validation and ablation
    study is defined for the revision; its results remain pending.
\end{itemize}
The present work does not claim field validation, statistical superiority,
unseen-speed robustness, or persistent-state streaming operation.

\section{Pantograph--Catenary System and Problem Formulation}
\label{sec:system}

\subsection{Pantograph--Catenary Dynamics}
The simulation model described in our previous work represents the catenary
using absolute nodal coordinate formulation (ANCF) elements and the
pantograph using a lumped mass--spring--damper model, coupled through a
penalty-based contact relation \cite{ale2026realtime}. Its validation against
benchmark and field data is reported in that work; the files used here are
simulation outputs, not field measurements.

For an ANCF element, let $\chi\in[0,l_0]$ be the reference element
coordinate and $\mathbf e(t)\in\mathbb R^{n_e}$ its generalized nodal
coordinates, including the position and slope degrees of freedom. The
shape matrix $\mathbf S(\chi)$ maps these coordinates to the spatial
position $\mathbf r(\chi,t)$. Position and deformation energy are
\begin{align}
    \mathbf r(\chi,t) &= \mathbf S(\chi)\mathbf e(t),
    \label{eq:ancf_position}\\
    U(\mathbf e) &= \frac{1}{2}\int_0^{l_0}
    \left(EA\varepsilon_l^2+EI\kappa^2\right)\,d\chi,
    \label{eq:strain_energy}
\end{align}
where $E$, $A$, and $I$ denote elastic modulus, cross-sectional area, and
second moment of area. Axial strain $\varepsilon_l$ and curvature $\kappa$
depend on $\mathbf e$ through the spatial derivatives of
$\mathbf r(\chi,t)$. The energy variation satisfies
$\delta U=(\nabla_{\mathbf e}U)^{\mathrm T}\delta\mathbf e$; hence the
elastic restoring term placed on the left-hand side of the equations of
motion is
\begin{equation}
    \mathbf Q(\mathbf e)=\nabla_{\mathbf e}U(\mathbf e).
    \label{eq:generalized_force}
\end{equation}
With this convention, the acceleration contains the negative energy
gradient as its restoring contribution.

After assembly and treatment of constrained degrees of freedom, write the
catenary and pantograph balances as
\begin{align}
    \mathbf M_C\ddot{\mathbf U}_C+
    \mathbf C_C\dot{\mathbf U}_C+\mathbf Q_C(\mathbf U_C)
    &=\mathbf F_C,
    \label{eq:catenary_dynamics}\\
    \mathbf M_P\ddot{\mathbf U}_P+
    \mathbf C_P\dot{\mathbf U}_P+\mathbf K_P\mathbf U_P
    &=\mathbf F_P.
    \label{eq:pantograph_dynamics}
\end{align}
Here $\mathbf U_C\in\mathbb R^{n_C}$ and
$\mathbf U_P\in\mathbb R^{n_P}$ contain the subsystem coordinates.
Each mass, damping, or stiffness matrix has the corresponding square
dimension. The catenary restoring vector $\mathbf Q_C$ is assembled from
element energy gradients; it need not be a constant matrix times
$\mathbf U_C$. The lumped pantograph uses a linear stiffness matrix
$\mathbf K_P$ in this representation.

The load vectors $\mathbf F_C$ and $\mathbf F_P$ include external loads
and the contact interaction. The contact law determines the scalar force
$F_c(t)$ from the coupled motion and applies opposite contact reactions to
the two subsystems at the moving interface. Thus $F_c$ is an outcome of
the coupled solution, not an independently prescribed input. We retain
this general description because the available learning data do not
provide the contact-law parameters or all internal simulator states.

Stack the coordinates as
$\mathbf U=[\mathbf U_C^{\mathrm T},\mathbf U_P^{\mathrm T}]^{\mathrm T}$,
with dimension $n=n_C+n_P$. The assembled system then has the
second-order form
\begin{equation}
    \mathbf M\ddot{\mathbf U}(t)+\mathbf C\dot{\mathbf U}(t)
    +\mathbf Q(\mathbf U(t))=\mathbf F(t),
    \label{eq:full_mechanical}
\end{equation}
where $\mathbf Q(\mathbf U)$ collects the restoring terms and
$\mathbf F(t)$ is shorthand for the loads evaluated along the coupled
trajectory. In particular, its contact contribution can depend on the
system state as well as time. If the mass matrix on the free coordinates
is nonsingular, multiplication by $\mathbf M^{-1}$ gives
\begin{equation}
\begin{aligned}
    \ddot{\mathbf U}(t)={}&-\mathbf M^{-1}\mathbf C\dot{\mathbf U}(t)
    -\mathbf M^{-1}\mathbf Q(\mathbf U(t))\\
    &+\mathbf M^{-1}\mathbf F(t).
\end{aligned}
\label{eq:full_acceleration}
\end{equation}
To expose the connection with a first-order ODE, introduce the physical
generalized velocity $\mathbf w=\dot{\mathbf U}$ and state
$\mathbf y=[\mathbf U^{\mathrm T},\mathbf w^{\mathrm T}]^{\mathrm T}$.
Then
\begin{equation}
    \frac{d\mathbf y}{dt}=
    \begin{bmatrix}
        \mathbf w\\
        -\mathbf M^{-1}\mathbf C\mathbf w
        -\mathbf M^{-1}\mathbf Q(\mathbf U)
        +\mathbf M^{-1}\mathbf F
    \end{bmatrix}.
    \label{eq:mechanical_first_order}
\end{equation}
The first block is kinematic; the second contains dissipation, restoring
behavior, and driving. PGND borrows this separation, rather than solving
the full finite-element equations. No calibrated mass, damping, stiffness,
or finite-element state is supplied to the learning model. There is no
known projection from $\mathbf y$ to its learned latent state, so the
learned operators must not be identified with the physical matrices in
\eqref{eq:full_mechanical}.

\subsection{Available Observations and Target}
Let the response and reference force at physical time $t_i$ be
\begin{equation}
    \mathbf x_i=[a(t_i),u(t_i)]^{\mathrm T},
    \qquad F_i=F_c(t_i),
    \label{eq:observation}
\end{equation}
where $a$ is panhead acceleration and $u$ is uplift. These are only two
response channels of a much higher-dimensional mechanical state. In
particular, they do not directly supply every velocity, deformation, or
contact variable required to invert \eqref{eq:full_mechanical}. The
learning problem is therefore a response-to-force prediction problem,
not direct algebraic recovery of the complete physical state.

For a recording with $N$ samples, an eligible target index $k$ has
$m\leq k\leq N-1$. Its input--target pair is
\begin{equation}
    \mathbf X_k=
    \begin{bmatrix}
        \mathbf x_{k-m}^{\mathrm T}\\[-1mm]
        \vdots\\[-1mm]
        \mathbf x_{k-1}^{\mathrm T}
    \end{bmatrix}\in\mathbb R^{m\times2},
    \qquad y_k=F_k.
    \label{eq:window_pair}
\end{equation}
Within a single window, we relabel these times locally as
$t_0<\cdots<t_{m-1}<t_m$. The $m$ observations
$\mathbf x_0,\ldots,\mathbf x_{m-1}$ are available when the prediction is
issued; $F_m$ is its target. Every compared predictor follows the mapping
\begin{equation}
    \hat F_m=\mathcal P_\Theta(\mathbf x_0,\ldots,\mathbf x_{m-1};
    t_0,\ldots,t_m).
    \label{eq:discrete_mapping}
\end{equation}
where $\Theta$ denotes the model parameters and the sampling schedule is
known. Neither $\mathbf x_m$ nor any force-history input is available at
inference. The reference force $F_m$ is used only as a training label or
evaluation reference. For uniform sampling with interval $\Delta t$, the
observed history spans $(m-1)\Delta t$ and the forecast horizon is
$\Delta t$; the interval from the first observation to the target is
$m\Delta t$. For example, $m=16$ and $\Delta t=1$~ms give a 15-ms observed
span followed by a 1-ms prediction interval.

Windows are constructed within one recording and one data partition.
Changing their stride changes the selected targets, not the spacing of
the $m$ observations inside each window. PGND and the recurrent baselines
reset their states for each window. Consequently, this is finite-history,
one-step-ahead estimation, not contemporaneous reconstruction or
unlimited-history streaming. A reconstructed force curve is assembled from
successive window predictions, not from a single persistent latent rollout.

\subsection{Training-Only Standardization}
Let $\mathcal R_{\mathrm{tr}}$ contain the retained training-row indices
$(r,k)$, with $r$ identifying the recording, and let
$N_{\mathrm{tr}}=|\mathcal R_{\mathrm{tr}}|$. For each scalar channel
$\xi\in\{a,u,F\}$, compute the population statistics
\begin{align}
    \mu_\xi&=\frac{1}{N_{\mathrm{tr}}}
        \sum_{(r,k)\in\mathcal R_{\mathrm{tr}}}\xi_{r,k},
        \label{eq:training_mean}\\
    s_\xi^2&=\frac{1}{N_{\mathrm{tr}}}
        \sum_{(r,k)\in\mathcal R_{\mathrm{tr}}}
        (\xi_{r,k}-\mu_\xi)^2.
        \label{eq:training_variance}
\end{align}
The implemented divisor is $\sigma_\xi=\max(s_\xi,10^{-12})$ to avoid
division by zero. Means and scales are fitted once on training rows,
before constructing overlapping windows, rather than separately per
window or with repeated rows counted multiple times. Define
$\bm\mu_x=[\mu_a,\mu_u]^{\mathrm T}$ and
$\bm\sigma_x=[\sigma_a,\sigma_u]^{\mathrm T}$. The model receives
\begin{equation}
    \tilde{\mathbf x}_i=(\mathbf x_i-\bm\mu_x)\oslash\bm\sigma_x,
    \qquad \tilde F_i=(F_i-\mu_F)/\sigma_F,
    \label{eq:standardization}
\end{equation}
where $\oslash$ is element-wise division. The same frozen statistics are
used for all validation and test windows and for every compared model.
Force predictions are transformed back to the original force scale before
reporting errors. Partitions are defined before scaling and window
construction, as detailed in Section~\ref{sec:experiments}; validation or
test targets do not enter \eqref{eq:training_mean} or
\eqref{eq:training_variance}.

\section{Observation-Conditioned Physics-Guided Neural Dynamics}
\label{sec:method}

The revision retains the structured dynamics and changes the readout.
We call the implemented latent-only control \emph{PGND-0}; the additive
readout below defines the proposed \emph{observation-conditioned PGND}.
Warm-up initialization is a separate extension, not a prerequisite for
testing the readout.

For each window, the model encodes the responses, initializes a latent
state, integrates its dynamics, and decodes the target force. Both readout
branches use the same encoded observations. The learned parameters are
shared across windows; the initial state and trajectory are window-specific.

\subsection{Latent Coordinates and Time Scale}
Define neural solver time and latent state by
\begin{equation}
\begin{aligned}
    \tau_i&=(t_i-t_0)/t_{\mathrm{unit}},\\
    \mathbf z(\tau)&=[\mathbf q(\tau)^{\mathrm T},
                      \mathbf v(\tau)^{\mathrm T}]^{\mathrm T}.
\end{aligned}
\label{eq:latent_state}
\end{equation}
The constant $t_{\mathrm{unit}}>0$ converts physical seconds to
dimensionless solver time, with $\tau_0=0$ in every window.
Here $\mathbf q,\mathbf v\in\mathbb R^d$ and
$\mathbf z\in\mathbb R^{2d}$. The configuration-like coordinates
$\mathbf q$ represent learned information about the response history;
their velocity-like partners are defined by
\begin{equation}
    \frac{d\mathbf q}{d\tau}=\mathbf v.
    \label{eq:latent_kinematics}
\end{equation}
All latent derivatives below are with respect to $\tau$, not physical
time $t$. The chain rule makes the distinction explicit:
\begin{equation}
    \frac{d\mathbf q}{dt}=\frac{\mathbf v}{t_{\mathrm{unit}}},
    \qquad
    \frac{d^2\mathbf q}{dt^2}=
    \frac{1}{t_{\mathrm{unit}}^2}\frac{d\mathbf v}{d\tau}.
    \label{eq:time_conversion}
\end{equation}
The reference implementation uses $t_{\mathrm{unit}}=10^{-3}$~s; a
physical sampling interval of 1~ms therefore corresponds to
$\Delta\tau=1$. The variable $\mathbf v$ is latent displacement per
solver-time unit, not measured panhead velocity, and neither latent block
is assigned a calibrated physical unit. A different time scale changes
the numerical vector field and regularization balance, so it must be held
fixed or explicitly accounted for in comparisons.

\subsection{Observation Encoding and Interpolation}
\label{subsec:observation}
A shared encoder maps each available standardized observation to
\begin{equation}
    \mathbf h_i=\mathcal E_\psi(\tilde{\mathbf x}_i)
    \in\mathbb R^{d_e},\qquad i=0,\ldots,m-1.
    \label{eq:observation_encoder}
\end{equation}
We use a two-layer multilayer perceptron (MLP), with hidden width $w$:
\begin{equation}
    \mathcal E_\psi(\mathbf s)=
    \mathbf W_{E,2}\tanh(\mathbf W_{E,1}\mathbf s+\mathbf b_{E,1})
    +\mathbf b_{E,2}.
    \label{eq:encoder_mlp}
\end{equation}
Here $\mathbf W_{E,1}\in\mathbb R^{w\times2}$,
$\mathbf W_{E,2}\in\mathbb R^{d_e\times w}$, and the biases have
dimensions $w$ and $d_e$. The nonlinearity acts element-wise, and the
output layer is linear. Encoding is pointwise: $\mathbf h_i$ depends only
on $\tilde{\mathbf x}_i$, with temporal processing performed by the ODE.
The initializer, residual, and readout MLPs below use the same
affine--tanh--affine pattern with different input and output dimensions.

An ODE solver evaluates the vector field between observation times, so
the discrete encodings must define a continuous input path.
The observation path is linearly interpolated within the available history:
\begin{equation}
\begin{aligned}
    \mathbf h(\tau)&=(1-\rho)\mathbf h_i+\rho\mathbf h_{i+1},\\
    \rho&=\frac{\tau-\tau_i}{\tau_{i+1}-\tau_i},
    \quad \tau\in[\tau_i,\tau_{i+1}],
\end{aligned}
\label{eq:linear_interpolation}
\end{equation}
for $i=0,\ldots,m-2$. Over the unobserved prediction interval, use
\begin{equation}
    \mathbf h(\tau)=\mathbf h_{m-1},
    \qquad \tau\in[\tau_{m-1},\tau_m].
    \label{eq:continuous_observation}
\end{equation}
Both interpolation endpoints in \eqref{eq:linear_interpolation} are
available when the prediction is issued at $t_{m-1}$. This is compatible
with the finite-history task, but does not establish zero-delay estimation
at every internal integration time. No interpolation uses $\mathbf x_m$.
Causality of the upstream filters that generated the spreadsheets is not
established by this rule.

Interpolation is performed on the encoded vectors, not on raw responses
before encoding. Since $\mathcal E_\psi$ is nonlinear, these operations
are generally not interchangeable. Holding the final encoding is an
explicit boundary rule for the one-step prediction interval; it does not
assume that the physical responses actually remain constant.

\subsection{Initialization and Optional Causal Warm-Up}
\label{subsec:initialization}
Choose an initialization index $K_0$ with $0\leq K_0\leq m-2$ and set
\begin{equation}
    \mathbf z(\tau_{K_0})=
    \mathcal E_{0,\omega}
    (\tilde{\mathbf x}_0,\ldots,\tilde{\mathbf x}_{K_0}).
    \label{eq:initial_state}
\end{equation}
More explicitly, the initializer receives the concatenated vector
\begin{equation}
    \mathbf s_{K_0}=
    [\tilde{\mathbf x}_0^{\mathrm T},\ldots,
     \tilde{\mathbf x}_{K_0}^{\mathrm T}]^{\mathrm T}
    \in\mathbb R^{2(K_0+1)}
    \label{eq:initializer_input}
\end{equation}
and maps it to $2d$ unconstrained coordinates, split into
$\mathbf q(\tau_{K_0})$ and $\mathbf v(\tau_{K_0})$.
Neither block is obtained by differentiating the measured uplift; the
initializer learns a state useful for predicting force under the chosen
latent dynamics.

PGND-0 and the first readout-only comparison use $K_0=0$, matching the
implemented single-observation initializer. The warm-up extension uses
$K_0+1$ initial observations, concatenated and mapped by a small MLP.
Integration begins at $\tau_{K_0}$, not at an earlier time whose initial
state would then depend on later observations. Warm-up observations are
part of the same $m$-sample budget, not additional data. No state is
carried across windows or recording boundaries.

Warm-up length remains to be selected using a predeclared validation-only
candidate set. Its benefit is untested, and the latent state is not assumed
to be uniquely recoverable from the available responses.

\subsection{Structured Vector Field}
\label{subsec:physics_guided}
Following the kinematic/acceleration split in
\eqref{eq:mechanical_first_order}, the latent dynamics are
\begin{equation}
\begin{aligned}
    \frac{d\mathbf q}{d\tau}&=\mathbf v,\\
    \frac{d\mathbf v}{d\tau}&=-\mathbf D_\alpha\mathbf v
    -\nabla_{\mathbf q}V_\beta(\mathbf q)\\
    &\quad+\mathbf B_\gamma\mathbf h(\tau)
    +\mathbf r_\theta(\tau),
\end{aligned}
\label{eq:final_pgnd}
\end{equation}
All four terms in the second equation are $d$-dimensional.
The matrices $\mathbf D_\alpha\in\mathbb R^{d\times d}$ and
$\mathbf B_\gamma\in\mathbb R^{d\times d_e}$ act on velocity and encoded
observations, respectively. The scalar function
$V_\beta:\mathbb R^d\to\mathbb R$ determines the restoring term.
The remaining nonlinear correction is
\begin{equation}
    \mathbf r_\theta(\tau)=f_\theta
    (\mathbf q(\tau),\mathbf v(\tau),\mathbf h(\tau),\tau).
    \label{eq:neural_residual}
\end{equation}
Here $f_\theta:\mathbb R^{2d+d_e+1}\to\mathbb R^d$ is an MLP acting
on the concatenation of its arguments; its final scalar input is the
window-relative time $\tau$. The term $-\mathbf D_\alpha\mathbf v$
damps the latent velocity, $-\nabla_{\mathbf q}V_\beta$ generates
restoring dynamics, and $\mathbf r_\theta$ supplies a flexible correction
to this restricted structure. The identity $d\mathbf q/d\tau=\mathbf v$
is imposed exactly at the vector-field level, not learned by a separate
network or encouraged through a loss penalty.

The term $\mathbf B_\gamma\mathbf h$ conditions the state on measured
responses. It is not a physical external force: acceleration and uplift
are responses of the coupled system rather than applied excitations.

\subsection{Damping, Potential, and Latent Energy Balance}
\label{subsec:energy}
Use positive-semidefinite damping and a quadratic latent potential:
\begin{align}
    \mathbf D_\alpha&=\mathbf L_D\mathbf L_D^{\mathrm T}\succeq0,
    \label{eq:damping_constraint}\\
    V_\beta(\mathbf q)&=\tfrac12\mathbf q^{\mathrm T}\mathbf K_\beta\mathbf q,
    \quad \mathbf K_\beta=\mathbf L_K\mathbf L_K^{\mathrm T}\succeq0,
    \label{eq:quadratic_potential}\\
    \nabla_{\mathbf q}V_\beta(\mathbf q)
    &=\tfrac12(\mathbf K_\beta+\mathbf K_\beta^{\mathrm T})\mathbf q
    =\mathbf K_\beta\mathbf q.
    \label{eq:linear_restoring}
\end{align}
The factors $\mathbf L_D,\mathbf L_K\in\mathbb R^{d\times d}$ are
unconstrained trainable matrices; they need not be triangular. For any
$\boldsymbol\xi\in\mathbb R^d$,
\begin{equation}
    \boldsymbol\xi^{\mathrm T}\mathbf D_\alpha\boldsymbol\xi
    =\|\mathbf L_D^{\mathrm T}\boldsymbol\xi\|_2^2\geq0,
    \label{eq:damping_psd_proof}
\end{equation}
and the same argument applies to $\mathbf K_\beta$. Consequently,
$V_\beta(\mathbf q)=\tfrac12\|\mathbf L_K^{\mathrm T}\mathbf q\|_2^2$
is bounded below by zero. Symmetry is what reduces its gradient to
$\mathbf K_\beta\mathbf q$ in \eqref{eq:linear_restoring}. The restoring
term therefore requires only a matrix--vector product, while gradients
can still propagate to both $\mathbf q$ and $\mathbf L_K$.

The quadratic potential preserves the implemented structural core; it
does not imply that the complete model is linear, because the encoder,
neural residual, initializer, and readouts are nonlinear. A nonlinear
potential is not introduced simultaneously with the new readout, so that
their effects are not confounded.

For the latent energy-like quantity
\begin{equation}
    \mathcal E(\tau)=\tfrac12\mathbf v^{\mathrm T}\mathbf v
    +V_\beta(\mathbf q),
    \label{eq:latent_energy}
\end{equation}
the chain rule gives
\begin{equation}
    \frac{d\mathcal E}{d\tau}
    =\mathbf v^{\mathrm T}\frac{d\mathbf v}{d\tau}
    +(\nabla_{\mathbf q}V_\beta)^{\mathrm T}
      \frac{d\mathbf q}{d\tau}.
    \label{eq:energy_chain_rule}
\end{equation}
Substitution of \eqref{eq:final_pgnd} yields
\begin{equation}
\begin{aligned}
    \frac{d\mathcal E}{d\tau}
    ={}&-\mathbf v^{\mathrm T}\mathbf D_\alpha\mathbf v
         -\mathbf v^{\mathrm T}\nabla_{\mathbf q}V_\beta\\
       &+\mathbf v^{\mathrm T}\mathbf B_\gamma\mathbf h
         +\mathbf v^{\mathrm T}\mathbf r_\theta
         +(\nabla_{\mathbf q}V_\beta)^{\mathrm T}\mathbf v.
\end{aligned}
    \label{eq:energy_substitution}
\end{equation}
The two potential-gradient terms are equal scalars with opposite signs
and cancel. Thus the full continuous-time balance becomes
\begin{equation}
    \frac{d\mathcal E}{d\tau}
    =-\mathbf v^{\mathrm T}\mathbf D_\alpha\mathbf v
    +\mathbf v^{\mathrm T}\mathbf B_\gamma\mathbf h
    +\mathbf v^{\mathrm T}\mathbf r_\theta.
    \label{eq:complete_energy_rate}
\end{equation}
With the observation-conditioned and residual terms removed,
\begin{equation}
    \frac{d\mathcal E}{d\tau}
    =-\|\mathbf L_D^{\mathrm T}\mathbf v\|_2^2\leq0.
    \label{eq:energy_dissipation}
\end{equation}
This establishes non-increasing energy only for the unforced latent
subsystem. The two driving contributions in
\eqref{eq:complete_energy_rate} can have either sign, and semidefinite
matrices may have null directions; decay of every coordinate is not
guaranteed. The identity neither proves stability of the driven model nor
identifies physical pantograph energy. The unit-mass convention is a
latent modeling choice, not a calibrated physical mass normalization.
Finally, this continuous-time identity does not guarantee the same
behavior after discretization; numerical integration error must be checked.

\subsection{Evolution and Force Readout}
\label{subsec:force_reconstruction}
Let $f_{\mathrm{PGND}}$ denote the complete $2d$-dimensional vector field
in \eqref{eq:final_pgnd}. Given the initialized state and the known encoded
path, its trajectory satisfies the integral equation
\begin{equation}
\begin{aligned}
    \mathbf z(\tau)={}&\mathbf z(\tau_{K_0})\\
    &+\int_{\tau_{K_0}}^{\tau}
    f_{\mathrm{PGND}}(\mathbf z(s),\mathbf h(s),s)\,ds,
    \quad \tau_{K_0}\leq\tau\leq\tau_m.
\end{aligned}
    \label{eq:node_integral}
\end{equation}
The solver evaluates $\mathbf h(s)$ by
\eqref{eq:linear_interpolation}--\eqref{eq:continuous_observation},
including at internal integration stages. Numerically, the target state is
\begin{equation}
\begin{aligned}
    \mathbf z_m=\operatorname{ODESolve}\big(&f_{\mathrm{PGND}},
    \mathbf z(\tau_{K_0}),\\
    &[\tau_{K_0},\tau_m],\mathbf h\big).
\end{aligned}
    \label{eq:ode_solver}
\end{equation}
The reference solver is fixed-step, fourth-order Runge--Kutta (RK4).
Its step size $\delta\tau$ is distinct from the observation spacing
$\tau_{i+1}-\tau_i$, although both are 1 in the reference configuration.
The last integration interval still evolves the state using its damping,
restoring, and residual terms; holding the input does not freeze the
state. With $K_0=0$ and $m=16$, the solve covers $\tau=0$ through
$\tau=16$, including the final interval from 15 to 16.

The original readout uses only this state,
\begin{equation}
    \hat{\tilde F}^{(0)}_m=\mathcal H_\phi(\mathbf z_m).
    \label{eq:latent_only_decoder}
\end{equation}
where $\mathcal H_\phi:\mathbb R^{2d}\to\mathbb R$ is a scalar decoder.
The proposed readout adds a direct observation-conditioned predictor
$\mathcal G_\eta:\mathbb R^{d_e}\to\mathbb R$:
\begin{equation}
    \hat{\tilde F}_m=
    \underbrace{\mathcal G_\eta(\mathbf h_{m-1})}_{\text{direct prediction}}
    +\underbrace{\mathcal H_\phi(\mathbf z_m)}_{\text{latent correction}},
    \label{eq:force_decoder}
\end{equation}
with physical force given by
\begin{equation}
    \hat F_m=\mu_F+\sigma_F\hat{\tilde F}_m.
    \label{eq:force_reconstruction}
\end{equation}
Both branches predict the same next-sample standardized force from
information available by $t_{m-1}$. The direct branch does not use a
measured force or the unavailable $\mathbf x_m$. Its purpose is to test
whether retaining a direct response--force mapping helps while the latent
path represents history-dependent effects.

The ODE residual $\mathbf r_\theta$ and the output correction
$\mathcal H_\phi$ have different roles: the former changes a
$d$-dimensional state derivative, whereas the latter contributes one
scalar to the predicted standardized force. Neither is a directly measured
physical force component. The two readout branches are learned jointly;
for any constant $c$, replacing $\mathcal G_\eta$ by
$\mathcal G_\eta+c$ and $\mathcal H_\phi$ by
$\mathcal H_\phi-c$ leaves the sum unchanged. Thus even their offsets are
not separately identifiable from force supervision, and the direct branch
could dominate the prediction.
Consequently, an accuracy gain over PGND-0 is insufficient evidence for the
value of physics guidance. The direct-only and unstructured-ODE controls
in Section~\ref{sec:planned} are essential.

For the initial revision, retain $d=d_e=16$ and hidden width $w=32$. The encoder
is a $2\!\to\!32\!\to\!16$ tanh MLP, the state decoder is
$32\!\to\!32\!\to\!1$, and the proposed direct branch is
$16\!\to\!32\!\to\!1$. Hidden layers use tanh and outputs are linear.
The residual MLP is $49\!\to\!32\!\to\!16$; its final layer is initialized
to zero in PGND-0. The candidate warm-up initializer has dimensions
$2(K_0+1)\!\to\!32\!\to\!32$. These are explicit candidate settings,
not established optima. The revised parameter count and initialization
details must be recorded when it is implemented.

\subsection{Learning Objective}
\label{subsec:objective}
Let $j=1,\ldots,B$ index the windows in a minibatch, each with its own
initialized state and observation path. Supervise only the final
next-sample force:
\begin{equation}
    \mathcal L_{\mathrm{rec}}=\frac1B\sum_{j=1}^{B}
    (\tilde F^{(j)}_m-\hat{\tilde F}^{(j)}_m)^2.
    \label{eq:reconstruction_loss}
\end{equation}
The reference label and prediction in each term refer to the same target
row, not to the last observed row. The intermediate states are not
supervised by a measured latent trajectory, and decoding them does not
create additional force-training targets. Because
$F_m-\hat F_m=\sigma_F(\tilde F_m-\hat{\tilde F}_m)$, the corresponding
minibatch force MSE in original units is
\begin{equation}
    \mathrm{MSE}_{\mathrm{force}}=\sigma_F^2\mathcal L_{\mathrm{rec}}.
    \label{eq:loss_units}
\end{equation}
This relation allows standardized training curves and force errors in
newtons to be distinguished without changing the prediction task.

Regularize the residual at evolved sample times, excluding the initial state:
\begin{equation}
    \mathcal L_{\mathrm{res}}=
    \frac{1}{B(m-K_0)}\sum_{j=1}^{B}\sum_{i=K_0+1}^{m}
    \|\mathbf r^{(j)}_\theta(\tau_i)\|_2^2.
    \label{eq:residual_loss}
\end{equation}
In each summand,
$\|\mathbf r_\theta\|_2^2=\sum_{\ell=1}^d r_{\theta,\ell}^2$;
the norm sums over latent components, rather than averaging over them.
The outer factors average over $B$ windows and the $m-K_0$ evolved
observation/target times. These are the sample-grid states, not the
internal RK4 stages. At $i=m$, the residual uses the held encoding
$\mathbf h_{m-1}$. This is the implemented discrete penalty, not a
time-weighted integral; irregular sampling would require a separately
declared weighting rule. The total objective is
\begin{equation}
    \mathcal L=\mathcal L_{\mathrm{rec}}
    +\lambda_{\mathrm{res}}\mathcal L_{\mathrm{res}},
    \label{eq:total_loss}
\end{equation}
with trainable parameters
\begin{equation}
    \Theta=\{\psi,\omega,\mathbf L_D,\mathbf L_K,
    \mathbf B_\gamma,\theta,\phi,\eta\}.
    \label{eq:trainable_parameters}
\end{equation}
All parameters are learned jointly using Adam and direct automatic
differentiation through the numerical ODE solution. To describe this
dependency, let $\mathbf z_{\ell+1}=\Phi_\ell(\mathbf z_\ell;\Theta)$
be one numerical solver step, with $\ell$ indexing integration steps
rather than observation samples. The map $\Phi_\ell$ includes the
parameter-dependent encoded input at each stage. If
$\mathbf S_\ell=\partial\mathbf z_\ell/\partial\Theta$, the chain rule is
\begin{equation}
    \mathbf S_{\ell+1}=
    \frac{\partial\Phi_\ell}{\partial\mathbf z_\ell}\mathbf S_\ell
    +\frac{\partial\Phi_\ell}{\partial\Theta}.
    \label{eq:solver_gradient}
\end{equation}
The initial sensitivity comes from the initializer. Reverse-mode autograd
applies the chain rule without forming full sensitivity matrices. It
reaches the encoder through both the input path and direct readout, and
the damping/stiffness factors through their matrix products. Gradients
are those of the discretized RK4 solve, not a continuous-adjoint solver.

The reference residual weight is $\lambda_{\mathrm{res}}=10^{-3}$ for
standardized force MSE. If the same objective were expressed using raw
force MSE while retaining the same latent penalty, its residual weight
would need to be multiplied by $\sigma_F^2$. Reusing the same numerical
weight would change the trade-off. Likewise, changing the latent dimension,
latent scale, or solver-time scale can change the penalty magnitude.
The penalty discourages large latent corrections but does not guarantee
that the mechanical terms dominate or that the energy rate is negative.
No physical force-balance or latent-trajectory supervision loss is added.

For fair model selection, validation uses the force-prediction term alone,
averaged over all validation targets, with the training scalers frozen.
The proposed best-validation rule is
\begin{equation}
    e^*=\underset{1\leq e\leq E_{\max}}{\operatorname{arg\,min}}
    \mathcal L_{\mathrm{rec,val}}(\Theta_e),
    \label{eq:validation_selection}
\end{equation}
where $e$ is the epoch and $E_{\max}$ the preset epoch budget. The residual
penalty is excluded from this selection metric because the recurrent
baselines do not have a corresponding latent term. Test targets are never
used for gradient updates or checkpoint selection. This rule is proposed
for the revised experiments, not claimed for the existing final-epoch runs.

\begin{algorithm}[t]
\caption{Training and Validation of Proposed PGND}
\label{alg:pgnd}
\KwIn{Row-disjoint training/validation data; $m$, $K_0$,
$\lambda_{\mathrm{res}}$; epoch budget $E_{\max}$}
\KwOut{Best-validation and final-epoch checkpoints}
Fit scaling on training rows only; freeze all statistics\;
Initialize $\Theta$; set best validation MSE to $+\infty$\;
\For{$e=1,\ldots,E_{\max}$}{
    \ForEach{minibatch of recording-local training windows}{
        Encode $\tilde{\mathbf x}_{0:m-1}$ by
        \eqref{eq:observation_encoder}\;
        Interpolate available history; hold $\mathbf h_{m-1}$
        over the final interval\;
        Initialize $\mathbf z(\tau_{K_0})$ using
        \eqref{eq:initial_state}\;
        Solve \eqref{eq:final_pgnd} to $\tau_m$\;
        Predict $\hat{\tilde F}_m$ with \eqref{eq:force_decoder}\;
        Compute \eqref{eq:total_loss}, backpropagate, and update $\Theta$\;
    }
    Compute validation MSE with parameters frozen\;
    Save improving checkpoint; append epoch history\;
}
Save final-epoch checkpoint separately\;
\end{algorithm}

Algorithm~\ref{alg:pgnd} specifies the proposed selection procedure.
The completed experiment saved final-epoch weights only; its results
are not retrospectively relabeled as best-validation results.

\section{Simulation Protocol and Provisional Results}
\label{sec:experiments}

\noindent\emph{Working-draft status: the result table and its interpretation
are provisional material for organizing the experiments. The retained
numbers are existing latent-only PGND-0 reference runs, not results for
the revised method. They must be replaced or expanded by the planned
controlled evaluation before final conclusions are written.}

\subsection{Dataset Provenance and Partitions}
The available collection contains 120 headerless, four-column XLS files:
eight case labels, speeds 300/350/380~km/h, and cutoff labels
20/50/100/150/200~Hz. The previous paper describes 96 groups at four cutoffs
\cite{ale2026realtime}; the imported collection additionally contains
100-Hz files. Three older XLSX files have a different layout and are not
combined with this collection. The case labels follow the supplied files;
their exact physical configuration mapping must be confirmed before
configuration-specific conclusions are drawn.

The completed comparison uses only the 20-Hz files: cases 1--6 at all
three speeds form the training/validation pool, and cases 7--8 at those
same speeds form the test set. Thus all three speeds occur in training;
this is not unseen-speed extrapolation. Each file contains distance,
acceleration, uplift, and contact force. Distance increments and filename
speeds imply approximately 1-ms sampling under constant-speed and unit
assumptions. Explicit timestamps and filter-generation code are not
available, so timing provenance and filter causality remain qualifications.

The common preprocessing procedure is:
\begin{enumerate}
    \item Read all numeric rows without interpreting the first row as a
    header. Trim 10\% of each recording at each end; reject non-finite data.
    \item Reserve the final 15\% of retained rows of each training-pool
    recording for validation. Entire test recordings remain separate.
    \item Fit feature-wise population means and standard deviations, and
    force scaling, using training rows only.
    \item Construct windows independently within each recording and row
    partition. No observation or target crosses a partition boundary.
\end{enumerate}
No additional filter is applied. Training and validation share recording
identities, but not rows; temporal correlation can remain. The supplied
cutoff variants affect both inputs and force targets and must not be
treated as input-noise variants with a common unfiltered ground truth.

With $m=16$, training target stride 16, and validation/test stride 1, the
completed experiment has 12,126 training, 33,983 validation, and 69,909 test
windows. Every input still contains 16 contiguous observations. The
training stride reduces the number of supervised targets; it is not a
claim of using every eligible training window.

The legacy loader discarded the first numeric row, joined recordings
before windowing, fitted a scaler without using its transformed values,
and partitioned overlapping windows. The revised handling above is common
to all models and changes the experimental protocol. Consequently, these
are internal comparisons, not direct comparisons with the earlier
paper's numerical scores.

\subsection{Implemented Models and Training Settings}
The PyTorch RNN, LSTM, and GRU baselines each contain two 32-unit recurrent
layers and a scalar linear output. CNN--GRU adds a 64-channel, width-1
ReLU convolution before two 32-unit GRU layers. The width-1 convolution
mixes channels at each time instant; it is not a temporal smoothing
filter. All models are unidirectional and reset at each window.
Their trainable parameter counts are 3,233, 12,833, 9,825, and 15,969,
respectively. PGND-0 contains 5,761 trainable parameters. The recurrent
architectures follow the legacy configurations, but framework and
initialization differences preclude bitwise Keras reproduction.

The completed comparison uses seed 0, 100 epochs, batch size 128, and Adam
with learning rate $10^{-3}$, $(\beta_1,\beta_2)=(0.9,0.999)$ and
$\epsilon=10^{-7}$. Models share the same seeded training-window order.
There is no early stopping, scheduler, clipping, weight decay, or dropout.
PGND-0 uses $K_0=0$, the latent-only readout
\eqref{eq:latent_only_decoder}, and $\lambda_{\mathrm{res}}=10^{-3}$.
Its damping and stiffness factors start at $0.1\mathbf I$, giving
$\mathbf D=\mathbf K=0.01\mathbf I$. All scores below use final-epoch
weights; equal epochs do not imply equal runtime or equal tuning budgets.

\subsection{Metrics and Provisional Reference Results}
For $n$ common test targets in newtons, define
\begin{align}
    \mathrm{MAE}&=\frac1n\sum_{j=1}^n|F_j-\hat F_j|,\\
    \mathrm{MSE}&=\frac1n\sum_{j=1}^n(F_j-\hat F_j)^2,\\
    \mathrm{RMSE}&=\sqrt{\mathrm{MSE}},\qquad
    R^2=1-\frac{\sum_j(F_j-\hat F_j)^2}
                  {\sum_j(F_j-\bar F)^2}.
\end{align}
MAE and RMSE have units N, MSE has units N$^2$, and $R^2$ is undefined
for constant targets. Table~\ref{tab:preliminary} reports the completed
single-seed comparison. It contains no results for the proposed additive
readout or warm-up extension.

\begin{table}[t]
\caption{Provisional 20-Hz reference results: 100 epochs, seed 0,
final-epoch weights, and 69,909 common test targets. PGND-0 denotes the
original latent-only model, not the proposed revision.}
\label{tab:preliminary}
\centering
\begin{tabular}{lrrr}
\hline
Model & MAE (N) & RMSE (N) & $R^2$\\
\hline
LSTM & 9.227 & 12.279 & 0.8911\\
GRU & 9.659 & 12.731 & 0.8829\\
RNN & 9.884 & 13.100 & 0.8760\\
CNN--GRU & 9.073 & 12.141 & 0.8935\\
PGND-0 & 9.385 & 12.591 & 0.8855\\
\hline
\end{tabular}
\end{table}

PGND-0 has approximately 3.7\% higher RMSE than CNN--GRU and higher RMSE
on each of the six test recordings. These results do not establish an
accuracy advantage from physics guidance. PGND-0's minimum validation MSE
was 0.06275 at epoch 91, whereas its final validation MSE was 0.07264,
both in standardized-force units. Since intermediate weights were not
saved, that history cannot establish the test score of the best-validation
checkpoint. The same limitation applies to the baseline histories.

PGND-0 training and per-epoch validation took 1,056.7~s in total (median
10.52~s per epoch) on an NVIDIA RTX 3080 Ti with PyTorch 2.8.0+cu128 and
\texttt{torchdiffeq} 0.2.5. This is not per-sample inference latency or a
real-time deployment benchmark. The completed recurrent baseline runs were
substantially faster. Execution optimizations cached invariant matrices
within each ODE solve and reduced transfers/synchronizations, without
changing the mathematical model. Training-time differences still need to
be reported alongside accuracy in future comparisons.

\section{Controlled Evaluation of the Proposed Revision}
\label{sec:planned}

The following experiments are planned, not completed. The current dataset
and measured PGND-0 results remain the development benchmark. The aim is
to distinguish architectural effects rather than assume that additional
complexity improves prediction.

\subsection{Staged Comparisons}
First, implement best-validation checkpoint saving for every model and
retain final weights separately, as in Algorithm~\ref{alg:pgnd}. Use
validation force MSE, not the PGND regularized training objective, for
selection. The first architectural comparison changes only the readout,
retaining $m=16$, $K_0=0$, the quadratic potential, solver, and reference
training settings.

The minimum controls are:
\begin{enumerate}
    \item \emph{PGND-0}: latent-only readout, reproducing the original
    architectural control under the new common checkpoint-selection rule.
    \item \emph{Observation-conditioned PGND}: the additive readout
    \eqref{eq:force_decoder} with otherwise unchanged dynamics.
    \item \emph{Direct-only predictor}: the same encoder and
    $\mathcal G_\eta(\mathbf h_{m-1})$, with no ODE or latent correction.
    This intentionally uses only the latest observation; it isolates the
    added prediction path, not a full-history baseline.
    \item \emph{Unstructured neural ODE}: replace the structured vector
    field with a learned vector field while retaining the state dimension,
    encoder, initializer, additive readout, input path, and solver. Report
    any parameter-count differences rather than claiming exact matching.
    \item \emph{CNN--GRU and LSTM}: retain the strongest measured recurrent
    references under the same data and selection protocol.
\end{enumerate}
Separate residual-weight and residual-removal ablations can then test the
regularization and learnable correction in the vector field. Removing
$\mathcal H_\phi$ from the readout is different from removing
$\mathbf r_\theta$ from the dynamics; these interventions must not share
an ambiguous \emph{no residual} label.

Next, test $m=64$ and, if justified by validation and cost, $m=128$ for
PGND and the recurrent references. The 16-sample reference represents
approximately 16~ms up to the prediction target; a 20-Hz oscillation has
a 50-ms period, but 20~Hz is a cutoff label, not an established dominant
signal frequency. Longer context is therefore a hypothesis to test, not
a known remedy. Compare lengths on common eligible target indices within
each partition so that changes in scored samples do not masquerade as
context benefits; report the resulting counts. Test the warm-up
initializer only after isolating the readout and window-length effects.
Its candidate lengths and initialization rules must be declared before
those runs.

\subsection{Selection, Repetitions, and Interpretation}
Use the same predefined seeds, for example 0--4, for all principal
comparisons, and report aggregate and per-recording metrics. Any tuning
of learning rate, residual weight, or architecture is restricted to
training/validation data, with candidate sets and search budgets reported
for each model. Seed variability measures optimization variability, not
diversity of physical operating conditions. The 69,909 overlapping test
windows are not independent experimental replicates; uncertainty analysis
must respect recording and temporal dependence.

The existing test scores have already informed the decision to revise the
method. Repeated use of these recordings is therefore development-stage
evidence, not an untouched confirmatory test. New held-out simulation runs
or a separately reserved, previously unexamined condition set will be
needed for final confirmation. Any future comparison across bandwidths
must group filtered versions of the same underlying simulation to avoid
mistaking related trajectories for independent holdouts.

A gain over PGND-0 alone supports the readout modification, not the physics
claim. Evidence for the mechanical structure requires comparison with the
unstructured ODE and appropriate ablations. If those comparisons reveal
no accuracy or robustness benefit, the conclusion and claimed contribution
must remain limited. Noise, missing-observation, and cross-condition
experiments may be added under predefined common perturbations and
information budgets, but no such advantage is claimed in this draft.

\section{Discussion and Limitations}
The direct readout and warm-up initializer address different hypothesized
limitations, without identifying latent coordinates with physical states.
Enforcing simulator force balance would require its parameters, relevant
states, and conventions; arbitrary physical losses should not be imposed
on unidentified latent coordinates.

Evidence remains limited to one seed, one bandwidth, known training speeds,
and reduced training targets. Simulation-model validation does not validate
the estimator on field measurements, and unknown upstream filtering
precludes a zero-latency sensing claim. PGND also remains more expensive
than the recurrent references. Accuracy, inference latency, memory, and
numerical convergence must be assessed jointly before deployment claims.

\section{Conclusion}
Observation-conditioned PGND combines a direct response-based prediction
with a structured latent dynamical correction; causal warm-up initialization
is a separate extension. The original latent-only model does not outperform
CNN--GRU in the completed 100-epoch simulation comparison. The proposed
revision remains unvalidated. Its controls will distinguish contributions
from observation access, temporal context, and mechanical structure.

\section{Data Availability}
The simulation outputs originate from the previous study
\cite{ale2026realtime}; original spreadsheets are preserved unchanged
locally. A public release, persistent identifier, and redistribution terms
have not been established. No access commitment is made in this draft.

\section{Code Availability}
The PyTorch implementation and saved results correspond to PGND-0.
The additive readout, warm-up, best-validation checkpointing, and planned
controls require subsequent implementation and testing. Run metadata record
file hashes, splits, scaling, settings, and seeds. A public archival code
release and version identifier remain to be supplied before submission.

\section*{Competing Interests}
The authors declare no competing interests.

\section*{Funding Declaration}
This work is supported by the Natural Science Foundation of China
(No.62361016) and the Fundamental Research Funds for the Central
Universities (2682024CX003).

\bibliographystyle{IEEEtran}
\bibliography{annot}
\end{document}
